
\documentclass[10pt,letterpaper]{article}

% ==========================================
% PACKAGES
% ==========================================
\usepackage[utf8]{inputenc}
\usepackage[T1]{fontenc}
\usepackage[english]{babel}
\usepackage[margin=2.2cm]{geometry}
\usepackage{graphicx}
\usepackage{amsmath,amssymb}
\usepackage{array}
\usepackage{longtable}
\usepackage{color}
\usepackage{fancyhdr}

% ==========================================
% COLORS
% ==========================================
\definecolor{datasheetblue}{rgb}{0.07,0.21,0.35}
\definecolor{lightblue}{rgb}{0.90,0.94,0.98}

% ==========================================
% DOCUMENT CONFIGURATION
% ==========================================
\setlength{\parindent}{0pt}
\setlength{\parskip}{6pt}
\setlength{\tabcolsep}{6pt}
\renewcommand{\arraystretch}{1.25}

\pagestyle{fancy}
\fancyhf{}
\fancyhead[L]{\small ANALOG ONE}
\fancyhead[R]{\small Class AB Audio Amplifier}
\fancyfoot[L]{\small Technical Datasheet}
\fancyfoot[R]{\small Page \thepage}
\renewcommand{\headrulewidth}{0.4pt}

\setlength{\headheight}{14pt}

% Section formatting without titlesec
\newcommand{\datasheetsection}[1]{
    \vspace{0.3cm}
    {\large\bfseries\color{datasheetblue}#1}
    \par\vspace{0.1cm}
    \hrule
    \vspace{0.15cm}
}

\newcommand{\datasheetsubsection}[1]{
    \vspace{0.15cm}
    {\normalsize\bfseries\color{datasheetblue}#1}
    \par
}

\begin{document}

% ==========================================
% TITLE
% ==========================================
\begin{center}

{\Large\bfseries\color{datasheetblue}
ANALOG ONE}

\vspace{0.3cm}

{\Huge\bfseries
CLASS AB AUDIO AMPLIFIER}

\vspace{0.2cm}

{\large
Discrete-Component Audio Power Amplifier}

\vspace{0.4cm}

\fcolorbox{datasheetblue}{lightblue}{
\parbox{0.85\textwidth}{
\centering
\Large\bfseries 20 W\\[2mm]
\normalsize
Class AB Audio Amplifier\\
Complementary MOSFET Output Stage\\
8 $\Omega$ Loudspeaker Application
}
}

\end{center}

\vspace{0.5cm}

\begin{tabular}{ll}
\hline
Document Type & Technical Datasheet\\
Product & Class AB Audio Amplifier\\
Developer & Analog One\\
Institution & Tecnologico de Costa Rica\\
Technology & Discrete BJT and MOSFET\\
\hline
\end{tabular}

% ==========================================
% GENERAL DESCRIPTION
% ==========================================
\datasheetsection{1. General Description}

The Analog One Class AB Audio Amplifier is a
discrete-component amplification system designed
to increase audio signal amplitude and deliver
power to an 8 $\Omega$ loudspeaker.

The amplifier consists of three functional stages:
a BJT preamplifier, a MOSFET voltage amplification
stage, and a complementary MOSFET Class AB
power output stage.

The preamplifier performs initial signal conditioning
and voltage amplification.

The intermediate stage provides the required
voltage swing, while the Class AB output stage
supplies the current necessary to drive the
loudspeaker.

The voltage and power amplification stages
operate with nominal symmetrical
$\pm24$ V power supplies.

The reference output power is 20 W.

A separate power-stage simulation recorded
22.8 W of average power delivered to an
8 $\Omega$ load under the specified test conditions.

% ==========================================
% FEATURES
% ==========================================
\datasheetsection{2. Key Features}

\begin{itemize}

\item Discrete-component audio amplifier.

\item Three-stage amplification architecture.

\item Common-emitter BJT preamplifier.

\item MOSFET voltage amplification stage.

\item Complementary Class AB output topology.

\item Nominal $\pm24$ V supply rails.

\item 20 W reference output power.

\item 8 $\Omega$ nominal loudspeaker load.

\item Adjustable MOSFET bias.

\item Source resistors for current stabilization.

\item Output RC stabilization network.

\item AC-coupled signal interfaces.

\end{itemize}

% ==========================================
% FUNCTIONAL DIAGRAM
% ==========================================
\datasheetsection{3. Functional Block Diagram}

\begin{figure}[htbp]
\centering

\fbox{
\begin{minipage}{0.75\textwidth}
\centering

\textbf{AUDIO INPUT}

\vspace{0.2cm}
$\downarrow$
\vspace{0.2cm}

\fcolorbox{datasheetblue}{lightblue}{
\parbox{0.65\textwidth}{
\centering
\textbf{BJT PREAMPLIFIER}\\
2N2222 Common-Emitter Stage
}
}

\vspace{0.2cm}

$\downarrow$

\vspace{0.2cm}

\fcolorbox{datasheetblue}{lightblue}{
\parbox{0.65\textwidth}{
\centering
\textbf{VOLTAGE AMPLIFIER}\\
IRF530 / IRF9540
}
}

\vspace{0.2cm}

$\downarrow$

\vspace{0.2cm}

\fcolorbox{datasheetblue}{lightblue}{
\parbox{0.65\textwidth}{
\centering
\textbf{CLASS AB POWER STAGE}\\
IRFZ44N / IRF4905
}
}

\vspace{0.2cm}

$\downarrow$

\vspace{0.2cm}

\fcolorbox{datasheetblue}{lightblue}{
\parbox{0.65\textwidth}{
\centering
\textbf{8 $\Omega$ LOUDSPEAKER}
}
}

\end{minipage}
}

\caption{Functional architecture of the complete amplifier.}

\end{figure}

The input signal passes through the
preamplifier, voltage amplification stage,
and Class AB output stage before
reaching the loudspeaker.

% ==========================================
% ELECTRICAL SPECIFICATIONS
% ==========================================
\datasheetsection{4. Electrical Specifications}

\begin{table}[htbp]
\centering
\small

\begin{tabular}{p{7.2cm}cc}
\hline
\textbf{Parameter} &
\textbf{Value} &
\textbf{Unit}\\
\hline

Positive supply voltage & +24 & V\\

Negative supply voltage & -24 & V\\

Nominal load resistance & 8 & $\Omega$\\

Reference output power & 20 & W\\

Simulated power-stage output & 22.8 & W\\

Reference RMS output voltage & 12.65 & V\\

Reference peak output voltage & 17.89 & V\\

Reference RMS output current & 1.58 & A\\

Reference peak output current & 2.24 & A\\

\hline
\end{tabular}

\caption{Main electrical characteristics.}
\end{table}

The reference electrical quantities correspond
to the calculated requirements for a
20 W sinusoidal output into an 8 $\Omega$ load.

The simulated 22.8 W result corresponds
to a separate power-stage test.

% ==========================================
% PREAMPLIFIER
% ==========================================
\datasheetsection{5. Preamplifier Stage}

\datasheetsubsection{5.1 Description}

The preamplifier uses a 2N2222 NPN
bipolar junction transistor operating
in a common-emitter configuration.

The emitter resistor provides
negative feedback and stabilizes
the operating point.

Input and output coupling capacitors
block unwanted DC components.

\datasheetsubsection{5.2 Electrical Parameters}

\begin{table}[htbp]
\centering
\small

\begin{tabular}{p{7cm}cc}
\hline
Parameter & Value & Unit\\
\hline

Transistor & 2N2222 & --\\

Supply voltage & 24 & V\\

Base voltage & 7.39 & V\\

Emitter voltage & 6.69 & V\\

Collector voltage & 9.81 & V\\

Quiescent collector current & 3.02 & mA\\

Collector-emitter voltage & 3.12 & V\\

Voltage gain & -2.13 & V/V\\

Transistor power dissipation & 9.42 & mW\\

Estimated junction temperature & 26.88 & $^\circ$C\\

\hline
\end{tabular}

\caption{Calculated preamplifier characteristics.}
\end{table}

% ==========================================
% VOLTAGE AMPLIFIER
% ==========================================
\datasheetsection{6. Voltage Amplification Stage}

\datasheetsubsection{6.1 Description}

The voltage amplification stage consists
of two IRF530 N-channel MOSFETs and
one IRF9540 P-channel MOSFET.

The IRF530 transistors provide
voltage amplification.

The IRF9540 operates as an active load.

The stage is designed to produce
approximately 20 V peak output
from a 2 V RMS input signal.

\datasheetsubsection{6.2 Electrical Characteristics}

\begin{table}[htbp]
\centering
\small

\begin{tabular}{p{7cm}cc}
\hline
Parameter & Value & Unit\\
\hline

Supply voltage & $\pm24$ & V\\

Input test voltage & 2 & V RMS\\

Target output voltage & 20 & V peak\\

Target voltage gain & 7.07 & V/V\\

Simulated gain at 1 kHz & 7.118 & V/V\\

Simulated gain at 20 Hz & 6.930 & V/V\\

Output fundamental at 1 kHz & 20.112 & V peak\\

Total harmonic distortion & 1.183 & \%\\

Test load resistance & 43 & k$\Omega$\\

\hline
\end{tabular}

\caption{Voltage amplification characteristics.}
\end{table}

The total harmonic distortion value
corresponds to the individual
voltage amplification stage,
measured through the ninth harmonic.

% ==========================================
% POWER STAGE
% ==========================================
\datasheetsection{7. Class AB Power Stage}

\datasheetsubsection{7.1 Description}

The power output stage uses
a complementary MOSFET arrangement.

The IRFZ44N N-channel MOSFET and
IRF4905 P-channel MOSFET reproduce
the positive and negative portions
of the amplified audio waveform.

The Class AB configuration reduces
crossover distortion while maintaining
moderate quiescent power consumption.

Adjustable bias networks establish
the transistor operating conditions.

\datasheetsubsection{7.2 Output Characteristics}

\begin{table}[htbp]
\centering
\small

\begin{tabular}{p{7cm}cc}
\hline
Parameter & Value & Unit\\
\hline

Nominal load resistance & 8 & $\Omega$\\

Reference output power & 20 & W\\

Simulated load power & 22.8 & W\\

Simulated RMS output voltage & 13.5 & V\\

Simulated RMS output current & 1.69 & A\\

Peak-to-peak output voltage & 38.2 & V\\

Peak-to-peak output current & 4.78 & A\\

Simulation frequency & 10 & kHz\\

\hline
\end{tabular}

\caption{Recorded Class AB power-stage results.}
\end{table}

% ==========================================
% POWER SUPPLY
% ==========================================
\datasheetsection{8. Power Supply Requirements}

The voltage and power amplification stages
operate from symmetrical DC supply rails.

\begin{table}[htbp]
\centering

\begin{tabular}{p{7cm}cc}
\hline
Supply & Voltage & Unit\\
\hline

Positive rail & +24 & V\\

Negative rail & -24 & V\\

Reference ground & 0 & V\\

\hline
\end{tabular}

\caption{Main supply requirements.}
\end{table}

The preamplifier uses a nominal
24 V single-supply configuration.

All stages must share a suitable
signal reference when interconnected.

% ==========================================
% THERMAL CHARACTERISTICS
% ==========================================
\datasheetsection{9. Thermal Characteristics}

The Class AB power stage dissipates
energy through the output MOSFETs.

Adequate thermal management is
necessary to maintain safe
operating temperatures.

\begin{table}[htbp]
\centering

\begin{tabular}{p{7cm}cc}
\hline
Device & Dissipation & Unit\\
\hline

IRFZ44N & 6.58 & W\\

IRF4905 & 6.63 & W\\

Combined dissipation & 13.21 & W\\

\hline
\end{tabular}

\caption{Recorded output-stage power dissipation.}
\end{table}

These values correspond to the
10 kHz reference simulation.

Suitable heat sinks are recommended
for the output MOSFETs.

Heat-sink selection must consider
ambient temperature, junction
temperature limits, and
thermal resistance.

% ==========================================
% COMPONENT SUMMARY
% ==========================================
\datasheetsection{10. Component Summary}

\small

\begin{longtable}{p{3cm}p{3.5cm}p{8cm}}

\hline
\textbf{Stage} &
\textbf{Component} &
\textbf{Function}\\
\hline

Preamplifier &
2N2222 &
Common-emitter voltage amplification\\

Voltage stage &
IRF530 &
N-channel voltage amplification\\

Voltage stage &
IRF9540 &
P-channel active load\\

Power stage &
IRFZ44N &
Positive output current delivery\\

Power stage &
IRF4905 &
Negative output current return\\

Power stage &
0.22 $\Omega$, 5 W &
MOSFET source resistors\\

Power stage &
100 $\Omega$ &
MOSFET gate resistors\\

Power stage &
10 k$\Omega$ &
Adjustable bias trimmers\\

Power stage &
4.7 $\mu$F &
Signal coupling capacitors\\

Output network &
10 $\Omega$, 2 W &
Zobel stabilization resistor\\

Output network &
100 nF &
Zobel stabilization capacitor\\

Load &
8 $\Omega$, 30 W &
Nominal loudspeaker rating\\

\hline

\caption{Main amplifier components.}

\end{longtable}

\normalsize

% ==========================================
% TYPICAL APPLICATION
% ==========================================
\datasheetsection{11. Typical Application}

The Analog One amplifier is intended
for audio signal amplification
using an external loudspeaker.

The typical operating sequence is:

\begin{enumerate}

\item Audio signal input.

\item BJT signal conditioning.

\item MOSFET voltage amplification.

\item Complementary Class AB current amplification.

\item Audio output to an 8 $\Omega$ loudspeaker.

\end{enumerate}

The implementation requires
proper supply connections,
bias adjustment, grounding,
and thermal management.

% ==========================================
% OPERATING NOTES
% ==========================================
\datasheetsection{12. Operating Recommendations}

\begin{itemize}

\item Verify power supply polarity before operation.

\item Adjust the quiescent bias without an input signal.

\item Verify the DC output voltage before connecting
the loudspeaker.

\item Use appropriate heat sinks for the
output MOSFETs.

\item Use correctly rated power resistors.

\item Keep high-current output connections short.

\item Verify the complete amplifier after
interconnecting all three stages.

\end{itemize}

% ==========================================
% PERFORMANCE VERIFICATION
% ==========================================
\datasheetsection{13. Performance Verification}

The documented electrical results
were obtained through circuit simulation.

The Class AB power-stage reference
simulation used a 10 kHz input signal
and an 8 $\Omega$ load.

The recorded average output power
was 22.8 W.

The voltage amplification stage
was separately evaluated using
a 2 V RMS input signal at 1 kHz.

The simulated voltage gain was
7.118 V/V, with total harmonic
distortion of 1.183\%.

These results correspond to
individual stage evaluations.

Further testing is necessary
to determine the complete
amplifier frequency response,
harmonic distortion, noise,
and continuous output-power rating.

% ==========================================
% DOCUMENT INFORMATION
% ==========================================
\datasheetsection{14. Document Information}

\begin{tabular}{ll}
\hline
Developer & Analog One\\

Product & Class AB Audio Amplifier\\

Institution & Tecnologico de Costa Rica\\

Department & School of Electronics Engineering\\

Document & Technical Datasheet\\

Status & Design and Simulation\\
\hline
\end{tabular}

\vspace{0.8cm}

\begin{center}

\textcolor{datasheetblue}{
\textbf{
ANALOG ONE -- AUDIO AMPLIFICATION SYSTEMS
}
}

\end{center}

\end{document}
